%=====================================================================
\section{Switch and controller: every channel between them}
\label{sec:channels}

\begin{figure}[H]\centering
\begin{tikzpicture}
  \node[swbox, text width=46mm, minimum height=86mm, align=left] (s) at (0,0) {\small\textbf{switch} (\code{simple\_switch\_grpc})\\[4pt]
    table \code{forward} (key: arrival port)\\[3pt]
    clone session 100 $\to$ CPU port 255\\[3pt]
    digest 1 \code{hs\_digest\_t}\\
    digest 2 \code{tls\_digest\_t}\\
    digest 3 \code{http\_digest\_t}\\
    digest 4 \code{ssh\_digest\_t}\\
    digest 5 \code{dhcp\_digest\_t}\\
    digest 6 \code{x509\_digest\_t}\\[3pt]
    header \code{packet\_in}:\\ offset (32), fromClient (8),\\ msgStart (8), msgType (8),\\ connId (32)\\[3pt]
    header \code{packet\_out}:\\ kind (8), index (8)};
  \node[ctbox, text width=50mm, minimum height=86mm, align=left] (c) at (11.6,0) {\small\textbf{controller} (\code{controller/main.py})\\[4pt]
    \code{install\_forwarding}: 1$\to$2, 2$\to$1\\[3pt]
    \code{add\_clone\_session(100, 255)}\\[3pt]
    \code{on\_handshake}: JA4T, JA4TS, JA4L\\
    \code{on\_tls}: JA4, JA4S\\
    \code{on\_http}: JA4H (per request)\\
    \code{on\_ssh}: JA4SSH (per window)\\
    \code{on\_dhcp}: JA4D / JA4D6 (per message)\\
    \code{on\_x509}: JA4X (per certificate)\\[3pt]
    \code{on\_packet\_in} $\to$ \code{Reassembler}\\[3pt]
    \code{packet\_out}: kind 1 joined message,\\ kind 2 SSH flush, kind 3 certificate};
  \draw[arr] ($(c.west)+(0,3.1)$) -- node[lbl]{table write (once, at start)} ($(s.east)+(0,3.1)$);
  \draw[arr] ($(c.west)+(0,2.3)$) -- node[lbl]{clone session (once)} ($(s.east)+(0,2.3)$);
  \draw[arr] ($(s.east)+(0,0.6)$) -- node[lbl]{\textbf{digests}: the values the switch read\\ (stream channel, acknowledged one by one)} ($(c.west)+(0,0.6)$);
  \draw[arr] ($(s.east)+(0,-1.6)$) -- node[lbl]{\textbf{packet-in}: a copied segment + its place\\ in the byte stream} ($(c.west)+(0,-1.6)$);
  \draw[arr] ($(c.west)+(0,-2.9)$) -- node[lbl]{\textbf{packet-out}: a joined message, a flush,\\ a certificate: read by the same parser} ($(s.east)+(0,-2.9)$);
\end{tikzpicture}
\caption{The interface. Everything goes over one P4Runtime gRPC connection (\code{controller/p4runtime\_client.py}):
the controller becomes primary (arbitration), loads \code{build/main.json} with its P4Info, writes the table entries
and the clone session, enables the six digests, and then only listens and answers.}
\end{figure}

\subsection{The six digests}
\begin{center}\footnotesize
\begin{tabularx}{\textwidth}{c l >{\raggedright\arraybackslash}X >{\raggedright\arraybackslash}X >{\raggedright\arraybackslash}p{30mm}}
\toprule
\textbf{id} & \textbf{struct} & \textbf{sent when} & \textbf{main fields} & \textbf{controller writes} \\
\midrule
1 & \code{hs\_digest\_t} & first SYN (JA4T), first SYN-ACK (JA4TS), the ACK after it (JA4L) & window, option kinds (up to 40), MSS, window scale, two latencies, two TTLs, connId & \code{handshake.py}: \code{64240\_2-4-8-1-3\_1460\_7}, \code{1029\_64} \\
2 & \code{tls\_digest\_t} & a complete hello, once per connection and direction & versions, SNI flag, ciphers ($\le$128), extensions ($\le$64), signature algorithms ($\le$64), ALPN first/last, overflow, connId & \code{tls.py}: sort, SHA-256 $\to$ \code{t13d1516h2\_\ldots} \\
3 & \code{http\_digest\_t} & a complete request header block (not twice for one sequence number) & method, version, cookie / referer flags, header count, language, names ($\le$512 B), cookie pairs ($\le$1,024 B), serial, connId, the request's sequence number (the controller lists a connection's requests in stream order) & \code{http.py} $\to$ \code{ge11cr04engb\_\ldots} \\
4 & \code{ssh\_digest\_t} & 8 lengths collected; window of 200 closed; flush; an earlier connection's last window on a reused port's SYN; an SSH banner on a port other than 22 & batch of 8 (direction + length), 4 counts, serial, connId, direction & \code{ssh.py}: modes $\to$ \code{c36s52\_c100s100\_c0s0} \\
5 & \code{dhcp\_digest\_t} & every DHCP / DHCPv6 message & type, size, i/d flags, options ($\le$64), requested codes ($\le$128), serial & \code{dhcp.py} $\to$ \code{disco0576nn\_\ldots} \\
6 & \code{x509\_digest\_t} & every certificate sent back by the controller & chain index, issuer / subject / extension OIDs ($\le$256 B each), connId & \code{x509.py}: hex, SHA-256 $\to$ \code{a373a9f83c6b\_\ldots} \\
\bottomrule
\end{tabularx}
\end{center}

\textbf{Digest settings} (\code{enable\_digest}): \code{max\_list\_size = 1} and \code{max\_timeout\_ns = 0} (send
each digest at once, alone), \code{ack\_timeout\_ns} = 1\,s. The controller acknowledges every list as it handles
it. BMv2 holds an unacknowledged digest in a filter and drops an identical one: the DHCP, HTTP and SSH digests carry
a serial number for that reason (Section~\ref{sec:dec-serial}). Every digest of a TCP connection carries the
connection's number (\code{connId}, ConnId in \code{main.p4}): the controller keeps one record per (endpoints,
number), so a new connection on reused ports gets its own record.

%---------------------------------------------------------------------
\subsection{A message split over segments, arriving out of order}
\label{sec:seq}

\begin{figure}[H]\centering
\begin{tikzpicture}[yscale=1.0]
  \node[hobox] (cl) at (0,0) {client};
  \node[swbox] (sw) at (5.5,0) {switch};
  \node[ctbox] (ct) at (11,0) {controller};
  \foreach \x in {0,5.5,11} \draw[thick, gray] (\x,-0.6) -- (\x,-16.6);
  \draw[arr] (0,-1.5) -- node[lbl]{SYN (ISN $c$), SYN-ACK, ACK} (5.5,-1.5);
  \node[swbox, right, text width=40mm, align=left] at (5.7,-2.4) {Stream: ISN known, direction not started};
  \draw[arr] (0,-3.5) -- node[lbl]{segment 2: bytes 1460..2099} (5.5,-3.5);
  \node[swbox, right, text width=48mm, align=left] at (5.7,-4.7) {rule (c): direction not started,\\ offset 1460 $<$ 16 KB $\to$ copy,\\ and forward it};
  \draw[arr] (5.5,-5.9) -- node[lbl]{packet-in: offset 1460} (11,-5.9);
  \node[ctbox, left, text width=36mm, align=left] at (10.9,-7.0) {keep bytes at 1460;\\ no message start yet};
  \draw[arr] (0,-8.3) -- node[lbl]{segment 1: bytes 0..1459 (hello start)} (5.5,-8.3);
  \node[swbox, right, text width=48mm, align=left] at (5.7,-9.5) {parser: 0x16, hello longer than\\ the segment $\to$ incomplete;\\ rule (a): pending, copy, msgStart};
  \draw[arr] (5.5,-10.8) -- node[lbl]{packet-in: offset 0, start, TLS} (11,-10.8);
  \node[ctbox, left, text width=40mm, align=left] at (10.9,-12.0) {bytes 0..2099 contiguous;\\ hello length satisfied $\to$\\ one record, the whole hello};
  \draw[arr] (11,-13.3) -- node[lbl]{packet-out kind 1: the joined hello} (5.5,-13.3);
  \node[swbox, right, text width=48mm, align=left] at (5.7,-14.5) {same parser walks it $\to$ Tls digest\\ (reassembled = 1); pending cleared;\\ the packet is dropped, not forwarded};
  \draw[arr] (5.5,-15.8) -- node[lbl]{digest 2: the hello's values} (11,-15.8);
\end{tikzpicture}
\caption{The out-of-order case end to end. The original segments are forwarded unchanged as they pass; only copies
go to the controller. The controller joins bytes by TLS framing only and never reads a field: the fingerprint is
still read by the switch's parser.}
\end{figure}

\begin{figure}[H]\centering
\begin{tikzpicture}
  \node[term] (s) at (0,0) {data segment, not from the controller};
  \node[dec] (a) at (0,-1.6) {(a) a message starts\\ here, continues later?};
  \node[swbox, text width=42mm] (ay) at (6.0,-1.6) {pending, pendingSeq := seq;\\ copy with msgStart, type TLS / HTTP};
  \node[dec] (d) at (0,-3.5) {(d) a complete TLS $\le$1.2\\ ServerHello?};
  \node[swbox, text width=42mm] (dy) at (6.0,-3.5) {pending; copy with msgStart,\\ type 3 (server flight: certificates)};
  \node[dec] (b) at (0,-5.4) {(b) pending and within\\ 16 KB of its start?};
  \node[swbox, text width=42mm] (by) at (6.0,-5.4) {copy (a later piece of the message)};
  \node[dec] (c) at (0,-7.3) {(c) direction not started,\\ offset $<$ 16 KB?};
  \node[swbox, text width=42mm] (cy) at (6.0,-7.3) {copy (a piece that may have arrived\\ before the message's first segment)};
  \node[dec] (e) at (0,-9.2) {(e) HTTP direction, ahead of\\ a gap, within 16 KB of it?};
  \node[swbox, text width=42mm] (ey) at (6.0,-9.2) {copy (a later piece of the next request\\ on a keep-alive connection)};
  \node[term] (no) at (0,-10.8) {no copy};
  \draw[arr] (s) -- (a); \draw[arr] (a) -- node[lbl, above]{yes} (ay);
  \draw[arr] (a) -- node[lbl, right]{no} (d); \draw[arr] (d) -- node[lbl, above]{yes} (dy);
  \draw[arr] (d) -- node[lbl, right]{no} (b); \draw[arr] (b) -- node[lbl, above]{yes} (by);
  \draw[arr] (b) -- node[lbl, right]{no} (c); \draw[arr] (c) -- node[lbl, above]{yes} (cy);
  \draw[arr] (c) -- node[lbl, right]{no} (e); \draw[arr] (e) -- node[lbl, above]{yes} (ey);
  \draw[arr] (e) -- node[lbl, right]{no} (no);
\end{tikzpicture}
\caption{\code{p4/stream.p4}: which segments the controller must see. Offsets are \code{seq} $-$ ISN $-$ 1 per
direction (ISN from the SYN / SYN-ACK; without them, learned from the first message start). A joined message (or
the certificates) coming back clears \code{pending}. For (e) the switch keeps, per direction, the \emph{edge} (the
sequence number after the furthest byte seen): a segment beyond it opens a gap at the edge; the gap closes when a
segment covering its first byte arrives, and is given up 16 KB on. Everything else is forwarded without a copy: the
controller sees only what a split message may need.}
\label{fig:copyrules}
\end{figure}

\begin{figure}[H]\centering
\begin{tikzpicture}
  \node[term] (in) at (0,0) {packet-in: (endpoints, connId, direction), seq, offset, msgStart, type, payload};
  \node[ctbox, text width=70mm] (keep) at (0,-1.2) {pieces[offset] := payload, keeping bytes already held\\ (duplicates harmless; on overlap, first-received bytes win)};
  \node[ctbox, text width=70mm] (st) at (0,-2.4) {msgStart: starts[offset] := type};
  \node[ctbox, text width=70mm] (loop) at (0,-3.6) {for every start not yet done: bytes := contiguous(pieces, start)\\ (up to the first gap, at most 16 KB + 9)};
  \node[dec] (k) at (0,-5.0) {type?};
  \node[ctbox, text width=34mm] (t) at (-4.6,-6.6) {TLS: strip record headers,\\ join fragments until the\\ handshake length is met $\to$\\ one record};
  \node[ctbox, text width=30mm] (h) at (0,-6.6) {HTTP: up to the first\\ empty line\\ (\code{\textbackslash r?\textbackslash n\textbackslash r?\textbackslash n})};
  \node[ctbox, text width=34mm] (f) at (4.6,-6.6) {server flight: handshake\\ messages until Certificate\\ (11) $\to$ its DER list;\\ ServerHelloDone first $\to$ []};
  \node[dec] (cpl) at (0,-8.4) {complete?};
  \node[term] (wait) at (4.6,-8.4) {wait for more bytes};
  \node[ctbox, text width=70mm] (out) at (0,-9.7) {mark done; return (seq of its first byte, type, message) $\to$ packet-out};
  \draw[arr] (in) -- (keep); \draw[arr] (keep) -- (st); \draw[arr] (st) -- (loop); \draw[arr] (loop) -- (k);
  \draw[arr] (k) -| (t); \draw[arr] (k) -- (h); \draw[arr] (k) -| (f);
  \draw[arr] (t) |- (cpl); \draw[arr] (h) -- (cpl); \draw[arr] (f) |- (cpl);
  \draw[arr] (cpl) -- node[lbl, above]{no} (wait); \draw[arr] (cpl) -- node[lbl, right]{yes} (out);
\end{tikzpicture}
\caption{\code{controller/reassembly.py}. It only knows framing (record headers, handshake lengths, the empty
line, DER lengths of the certificate list) and never interprets a fingerprint field. A ``start'' flagged on a short
segment that only looked like a message never completes and is ignored. Bounded: at most 32 KB per direction (lowest
offsets dropped first), a direction idle for 60\,s is dropped, at most 100,000 directions. The bytes of a joined
message are not dropped at once: a false start in the middle of a request can be ``joined'' before the real one,
whose first segment arrives last, and that one still needs them.}
\end{figure}

\subsection{TLS 1.2 certificates (JA4X) and the SSH flush}

\begin{figure}[H]\centering
\begin{tikzpicture}[yscale=1.0]
  \node[hobox] (sv) at (0,0) {server};
  \node[swbox] (sw) at (5.5,0) {switch};
  \node[ctbox] (ct) at (11,0) {controller};
  \foreach \x in {0,5.5,11} \draw[thick, gray] (\x,-0.6) -- (\x,-11.4);
  \draw[arr] (0,-1.5) -- node[lbl]{ServerHello (TLS 1.2) + Certificate\ldots} (5.5,-1.5);
  \node[swbox, right, text width=46mm, align=left] at (5.7,-2.6) {JA4S digest; rule (d): copy the\\ flight from the ServerHello on};
  \draw[arr] (0,-3.8) -- node[lbl]{the rest of the flight (several segments)} (5.5,-3.8);
  \draw[arr] (5.5,-4.8) -- node[lbl]{packet-in (rule b)} (11,-4.8);
  \node[ctbox, left, text width=40mm, align=left] at (10.9,-6.0) {Certificate message complete:\\ cut it into its DER certificates};
  \draw[arr] (11,-7.3) -- node[lbl]{packet-out kind 3, index $i$ (one per certificate)} (5.5,-7.3);
  \node[swbox, right, text width=46mm, align=left] at (5.7,-8.5) {\code{parser\_x509.p4} walks the DER;\\ X509 digest with index $i$};
  \draw[arr] (5.5,-9.8) -- node[lbl]{digest 6 per certificate} (11,-9.8);
  \node[ctbox, left, text width=40mm, align=left] at (10.9,-10.8) {JA4X list in chain order};
\end{tikzpicture}
\caption{JA4X. Certificates are long (1--2 KB each) and always span segments, so they always take the join path.}
\end{figure}

\begin{figure}[H]\centering
\begin{tikzpicture}[yscale=1.0]
  \node[swbox] (sw) at (0,0) {switch};
  \node[ctbox] (ct) at (8,0) {controller};
  \foreach \x in {0,8} \draw[thick, gray] (\x,-0.6) -- (\x,-7.4);
  \node[swbox, right, text width=56mm, align=left] at (0.2,-1.4) {every SSH segment: counts in registers; every 8\\ lengths a BATCH digest; every 200 packets a WINDOW digest};
  \draw[arr] (0,-2.8) -- node[lbl]{digests 4 (kind 2, 3)} (8,-2.8);
  \node[ctbox, left, text width=46mm, align=left] at (7.8,-3.9) {traffic over: \code{finish()}};
  \draw[arr] (8,-5.0) -- node[lbl]{packet-out kind 2 (flush) with the connection's endpoints} (0,-5.0);
  \draw[arr] (0,-6.4) -- node[lbl]{digest 4 kind 4 (FINAL): the partial window} (8,-6.4);
\end{tikzpicture}
\caption{The last JA4SSH window. FoxIO reports the partial last window at the end of its capture; P4Runtime on
BMv2 cannot read registers (UNIMPLEMENTED), so the switch is asked to report it itself.}
\end{figure}
